\documentclass[11pt,a4paper]{article}
\usepackage[utf8]{inputenc}
\usepackage[spanish,es-tabla]{babel}
\usepackage{amsmath,amsfonts,amssymb}
\usepackage{geometry}
\usepackage{graphicx}
\usepackage{xcolor}
\usepackage{enumitem}
\usepackage{tikz}
\usetikzlibrary{positioning,shapes,arrows.meta,calc}

\geometry{margin=2.5cm}

\title{\textbf{Solución Examen Final 2020\\TAF ISC\\IMT Atlantique}}
\author{}
\date{}

\begin{document}

\maketitle

\section*{Ejercicio 1 - OFDM para modem acústico submarino}

\textbf{Contexto del problema:}

Sistema OFDM con las siguientes características:
\begin{itemize}
    \item Número de subportadoras: $N = 32$ alrededor de 12 kHz
    \item Modulación: PSK-8 en cada subportadora
    \item Código corrector de rendimiento: $r_c = \frac{2}{3}$
    \item Retardo máximo del canal: $\tau_{max} = 40$ ms
    \item Tiempo de coherencia del canal: $T_{coh} = 250$ ms
    \item Duración de un símbolo OFDM: $T_{OFDM} = 250$ ms
    \item Velocidad de propagación del sonido: $v = 1500$ m/s
\end{itemize}

\subsection*{Pregunta 1a: Distancia entre subportadoras}

\textbf{Enunciado:} Dar el valor de la distancia entre subportadoras.

\vspace{0.3cm}
\textbf{Solución paso a paso:}

\textbf{Paso 1: Relación fundamental en OFDM}

\textit{¿Por qué esta relación?} En OFDM, las subportadoras deben ser ortogonales entre sí para evitar interferencia. Esta ortogonalidad se logra cuando la separación en frecuencia es el inverso de la duración del símbolo OFDM.

La distancia entre subportadoras está dada por:
\[
\Delta f = \frac{1}{T_{OFDM}}
\]

\textit{Concepto clave:} Esta relación garantiza que durante un periodo de símbolo OFDM, cada subportadora complete un número entero de ciclos más que sus vecinas, asegurando ortogonalidad:
\[
\int_0^{T_{OFDM}} e^{i2\pi(f_k - f_j)t} dt = 0 \quad \text{para } k \neq j
\]
cuando $f_k - f_j = n\Delta f$ con $n$ entero.

\textbf{Paso 2: Calcular la distancia entre subportadoras}

Con $T_{OFDM} = 250$ ms:
\[
\Delta f = \frac{1}{250 \times 10^{-3}} = \frac{1}{0.25} = 4 \text{ Hz}
\]

\textbf{Respuesta:}

\[
\boxed{\Delta f = 4 \text{ Hz}}
\]

\textbf{Nota sobre eficiencia espectral:}
La separación de 4 Hz permite que las subportadoras se solapen en frecuencia sin causar interferencia gracias a la ortogonalidad. Sin ortogonalidad, habría ICI (Inter-Carrier Interference).

\textbf{Nota importante sobre notación:} En OFDM, $T_{OFDM}$ normalmente se refiere al \textbf{tiempo útil} (duración de la IFFT), no al símbolo completo con CP. Esta es la convención estándar porque:
\begin{itemize}
    \item La ortogonalidad depende solo del tiempo útil: $\Delta f = 1/T_{OFDM}$
    \item El CP se añade después para combatir ISI, pero no afecta la ortogonalidad
    \item Permite calcular directamente la separación entre subportadoras
\end{itemize}

Por tanto, $T_{OFDM} = 250$ ms es el \textbf{tiempo útil} donde se transmiten las subportadoras ortogonales.

\subsection*{Pregunta 1b: Duración del prefijo cíclico}

\textbf{Enunciado:} Dar la duración del prefijo cíclico.

\vspace{0.3cm}
\textbf{Solución:}

\textbf{Principio:}

El prefijo cíclico debe ser al menos igual al retardo máximo del canal para:
\begin{itemize}
    \item Eliminar la ISI entre símbolos OFDM consecutivos
    \item Mantener la ortogonalidad entre subportadoras
\end{itemize}

\textbf{Condición:}
\[
T_{CP} \geq \tau_{max}
\]

\textbf{Valor mínimo:}
\[
T_{CP,min} = \tau_{max} = 40 \text{ ms}
\]

En la práctica, a menudo se toma un margen de seguridad, pero basándonos en los datos del problema:

\textbf{Respuesta:}

\[
\boxed{T_{CP} = 40 \text{ ms}}
\]

\textbf{Verificación de tiempos:}

En OFDM, el símbolo completo transmitido tiene dos partes:
\begin{itemize}
    \item Parte útil (IFFT): $T_{OFDM} = 250$ ms (donde $\Delta f = 1/T_{OFDM}$)
    \item Prefijo cíclico: $T_{CP} = 40$ ms (añadido para combatir ISI)
\end{itemize}

\textbf{Duración total del símbolo transmitido:}
\[
T_{total} = T_{OFDM} + T_{CP} = 250 + 40 = 290 \text{ ms}
\]

\textbf{Nota conceptual:} El CP es una \textbf{copia del final} del símbolo útil que se añade al principio. No afecta la ortogonalidad (que depende solo de $T_{OFDM}$), pero sí aumenta el tiempo total de transmisión y reduce la eficiencia espectral.

\subsection*{Pregunta 1c: Número de bits por símbolo OFDM}

\textbf{Enunciado:} Dar el número de bits por símbolo OFDM.

\vspace{0.3cm}
\textbf{Solución paso a paso:}

\textbf{Paso 1: Bits por símbolo PSK-8}

PSK-8 tiene $M = 8$ símbolos, por tanto:
\[
m = \log_2(8) = 3 \text{ bits por símbolo}
\]

\textbf{Paso 2: Número total de subportadoras}

\[
N = 32 \text{ subportadoras}
\]

\textbf{Paso 3: Bits codificados por símbolo OFDM}

Si todas las subportadoras transmiten datos:
\[
\text{Bits codificados} = N \times m = 32 \times 3 = 96 \text{ bits}
\]

\textbf{Paso 4: Efecto del código corrector}

Con un código de rendimiento $r_c = \frac{2}{3}$:

\textit{Interpretación:} Por cada 2 bits de información, se generan 3 bits codificados.
\[
\text{Bits de información} = 96 \times \frac{2}{3} = 64 \text{ bits}
\]

\textbf{Respuesta (sin pilotos):}

\[
\boxed{\text{Bits de información por símbolo OFDM} = 64 \text{ bits}}
\]

\textbf{Nota:} Este valor puede reducirse si algunas subportadoras se dedican a pilotos (como veremos en las preguntas siguientes).

\subsection*{Pregunta 1d: Banda de frecuencias utilizada}

\textbf{Enunciado:} Dar la banda de frecuencias utilizada.

\vspace{0.3cm}
\textbf{Solución paso a paso:}

\textbf{Paso 1: Ancho de banda OFDM}

El ancho de banda ocupado por las $N$ subportadoras es:
\[
B = N \times \Delta f
\]

\textbf{Paso 2: Calcular el ancho de banda}

\[
B = 32 \times 4 = 128 \text{ Hz}
\]

\textbf{Paso 3: Frecuencia central}

El problema indica que las subportadoras están "alrededor de 12 kHz", por tanto:
\[
f_c = 12 \text{ kHz}
\]

\textbf{Paso 4: Banda de frecuencias}

Si las subportadoras están centradas en 12 kHz:
\[
f_{min} = f_c - \frac{B}{2} = 12000 - \frac{128}{2} = 12000 - 64 = 11936 \text{ Hz}
\]
\[
f_{max} = f_c + \frac{B}{2} = 12000 + \frac{128}{2} = 12000 + 64 = 12064 \text{ Hz}
\]

\textbf{Respuesta:}

\[
\boxed{B = 128 \text{ Hz}}
\]

\[
\boxed{\text{Banda: } [11.936 \text{ kHz}, 12.064 \text{ kHz}] \text{ o bien } 12 \text{ kHz} \pm 64 \text{ Hz}}
\]

\textbf{Observación importante:}

El ancho de banda de 128 Hz es muy estrecho, típico de las comunicaciones acústicas submarinas debido a:
\begin{itemize}
    \item La baja velocidad de propagación (1500 m/s vs $3\times10^8$ m/s para ondas EM)
    \item El largo retardo del canal (40 ms)
    \item Las limitaciones del medio submarino
\end{itemize}

\subsection*{Pregunta 1e: Débito binario (información útil)}

\textbf{Enunciado:} Dar el débito binario de información útil alcanzado.

\vspace{0.3cm}
\textbf{Solución paso a paso:}

\textbf{Paso 1: Bits de información por símbolo OFDM}

De la pregunta 1c (sin considerar pilotos aún):
\[
\text{Bits/OFDM} = 64 \text{ bits}
\]

\textbf{Paso 2: Duración total por símbolo}

El tiempo total incluye el CP:
\[
T_{total} = T_{OFDM} + T_{CP} = 250 + 40 = 290 \text{ ms} = 0.29 \text{ s}
\]

\textbf{Paso 3: Débito binario útil}

La tasa se calcula con el tiempo total de transmisión:
\[
R_b = \frac{\text{Bits de información}}{T_{total}} = \frac{64}{0.29} \approx 220.69 \text{ bits/s}
\]

\textbf{Respuesta:}

\[
\boxed{R_b \approx 220.7 \text{ bits/s}}
\]

\textbf{Observación:} Este débito es bajo debido a:
\begin{itemize}
    \item La larga duración del símbolo completo (290 ms = 250 ms útil + 40 ms CP)
    \item El overhead del prefijo cíclico: $\frac{40}{290} \times 100\% \approx 13.8\%$
    \item El código corrector de rendimiento 2/3 (reduce bits útiles a 2/3 de los transmitidos)
    \item El medio acústico submarino (velocidad del sonido 1500 m/s, mucho menor que las ondas EM)
\end{itemize}

\textbf{Eficiencia del CP:}
\[
\eta_{CP} = \frac{T_{OFDM}}{T_{total}} = \frac{250}{290} \approx 0.862 = 86.2\%
\]

\subsection*{Pregunta 2: Configuración óptima de pilotos en frecuencia}

\textbf{Enunciado:} Se desea estimar el canal insertando símbolos pilotos. Ciertas subportadoras están dedicadas a transmitir estos pilotos. ¿Cuál de las configuraciones siguientes es la más adaptada garantizando el mejor débito binario útil?

\begin{itemize}
    \item[(a)] 1 piloto cada 3 subportadoras: $\{f_0, f_3, f_6, ...\}$
    \item[(b)] 1 piloto cada 7 subportadoras: $\{f_0, f_7, f_{14}, ...\}$
    \item[(c)] 1 piloto cada 10 subportadoras: $\{f_0, f_{10}, f_{20}, ...\}$
\end{itemize}

\vspace{0.3cm}
\textbf{Solución paso a paso:}

\textbf{Paso 1: Criterio de Nyquist en frecuencia}

\textit{¿Por qué necesitamos pilotos?} Los pilotos son símbolos conocidos que se insertan para estimar cómo el canal afecta a nuestra señal. Sin ellos, no podríamos ecualizar correctamente.

\textit{Teorema de muestreo aplicado al canal:} Para estimar un canal con retardo máximo $\tau_{max}$, necesitamos "muestrear" en frecuencia lo suficientemente rápido. El criterio de Nyquist en frecuencia nos dice:
\[
\Delta f_{piloto} \leq \frac{1}{2\tau_{max}}
\]

\textit{Interpretación física:} Si el canal tiene ecos que duran hasta $\tau_{max}$, en frecuencia el canal varía en un ancho de aproximadamente $1/\tau_{max}$. Para capturar estas variaciones, necesitamos pilotos separados por la mitad de esta distancia (Nyquist).

Donde $\tau_{max}$ es el delay spread del canal.

\textbf{Paso 2: Calcular el espaciamiento máximo permitido}

\[
\Delta f_{piloto,max} = \frac{1}{2 \times 40 \times 10^{-3}} = \frac{1}{0.08} = 12.5 \text{ Hz}
\]

\textbf{Paso 3: Evaluar cada opción}

\textbf{Opción (a):} 1 piloto cada 3 subportadoras
\begin{itemize}
    \item Número de pilotos: $\lceil 32/3 \rceil = 11$ pilotos
    \item Espaciamiento: $\Delta f_a = 3 \times 4 = 12$ Hz
    \item $\Delta f_a = 12 \text{ Hz} < 12.5$ Hz $\Rightarrow$ \textbf{CUMPLE el criterio} \checkmark
\end{itemize}

\textbf{Opción (b):} 1 piloto cada 7 subportadoras
\begin{itemize}
    \item Número de pilotos: $\lceil 32/7 \rceil = 5$ pilotos
    \item Espaciamiento: $\Delta f_b = 7 \times 4 = 28$ Hz
    \item $\Delta f_b = 28 \text{ Hz} > 12.5$ Hz $\Rightarrow$ \textbf{NO cumple el criterio}
\end{itemize}

\textbf{Opción (c):} 1 piloto cada 10 subportadoras
\begin{itemize}
    \item Número de pilotos: $\lceil 32/10 \rceil = 4$ pilotos
    \item Espaciamiento: $\Delta f_c = 10 \times 4 = 40$ Hz
    \item $\Delta f_c = 40 \text{ Hz} > 12.5$ Hz $\Rightarrow$ \textbf{NO cumple el criterio}
\end{itemize}

\textbf{Paso 4: Conclusión}

\textbf{Análisis comparativo:}
\begin{itemize}
    \item Opción (a): 12 Hz < 12.5 Hz - \textcolor{green}{CUMPLE el criterio de Nyquist}
    \item Opción (b): 28 Hz (124\% por encima del criterio)
    \item Opción (c): 40 Hz (220\% por encima del criterio)
\end{itemize}

\textit{Cálculo del espaciado máximo:}
\[
N_{max} = \left\lfloor \frac{\Delta f_{piloto,max}}{\Delta f} \right\rfloor = \left\lfloor \frac{12.5}{4} \right\rfloor = \lfloor 3.125 \rfloor = 3
\]

Podemos poner pilotos separados hasta 3 subportadoras y cumplir Nyquist.

\textbf{Respuesta:}

\[
\boxed{\text{Opción (a): 1 piloto cada 3 subportadoras}}
\]

\textbf{Justificación:}
\begin{itemize}
    \item \textbf{Cumple el criterio de Nyquist}: 12 Hz < 12.5 Hz
    \item Proporciona estimación robusta del canal (11 pilotos)
    \item Mejor compromiso entre estimación del canal y débito útil
    \item Subportadoras de datos: $32 - 11 = 21$
    \item Permite reconstruir $H(f)$ sin aliasing
\end{itemize}

\textbf{¿Por qué no opciones (b) o (c)?}
\begin{itemize}
    \item Violan Nyquist: causarían aliasing en la estimación del canal
    \item Mayor débito pero estimación pobre del canal
    \item Probable aumento de BER por mala ecualización
\end{itemize}

\subsection*{Pregunta 3: Repetición de pilotos en el tiempo}

\textbf{Enunciado:} Se repite el patrón de pilotos en el tiempo. Encontrar la buena repetición entre:

\begin{itemize}
    \item[(a)] Repetir en todos los símbolos OFDM
    \item[(b)] Repetir cada 2 símbolos OFDM
    \item[(c)] Repetir cada 5 símbolos OFDM
\end{itemize}

\vspace{0.3cm}
\textbf{Solución paso a paso:}

\textbf{Paso 1: Criterio de Nyquist en tiempo}

\textit{El canal cambia con el tiempo:} En comunicaciones móviles o acústicas submarinas, el canal no es constante. El movimiento del receptor, las corrientes marinas, etc., hacen que la respuesta del canal varíe.

\textit{Tiempo de coherencia $T_{coh}$:} Es el tiempo durante el cual el canal permanece aproximadamente constante. Después de este tiempo, nuestra estimación del canal queda obsoleta.

\textit{Criterio de muestreo temporal:} Para seguir las variaciones temporales del canal, debemos actualizar nuestra estimación (con pilotos) al menos cada:
\[
T_{piloto} \leq \frac{T_{coh}}{2}
\]

\textit{Analogía:} Si quieres seguir a alguien que se mueve, necesitas mirarlo frecuentemente. Si solo miras cada mucho tiempo, perderás su posición. Aquí $T_{coh}$ es el tiempo que tarda el canal en "moverse significativamente".

\textbf{Paso 2: Calcular el período máximo permitido}

\[
T_{piloto,max} = \frac{T_{coh}}{2} = \frac{250 \text{ ms}}{2} = 125 \text{ ms}
\]

\textbf{Paso 3: Calcular cuántos símbolos OFDM caben en el período máximo}

\[
N_{OFDM,max} = \left\lfloor \frac{T_{piloto,max}}{T_{OFDM}} \right\rfloor = \left\lfloor \frac{125}{250} \right\rfloor = \lfloor 0.5 \rfloor = 0
\]

\textbf{Interpretación crítica:} $N_{OFDM,max} = 0$ significa que no podemos esperar ni siquiera 1 símbolo OFDM completo. Por tanto, necesitamos pilotos en \textbf{cada símbolo OFDM}.

\textbf{Paso 4: Evaluar cada opción}

\textbf{Opción (a):} Pilotos en todos los símbolos OFDM
\begin{itemize}
    \item Período: $T_a = T_{OFDM} = 250$ ms
    \item Relación: $\frac{T_a}{T_{piloto,max}} = \frac{250}{125} = 2$
    \item Actualiza cada $T_{coh}$ (caso límite aceptable)
    \item \textbf{Es la única opción viable}
\end{itemize}

\textbf{Opción (b):} Pilotos cada 2 símbolos OFDM
\begin{itemize}
    \item Período: $T_b = 2 \times 250 = 500$ ms
    \item Excede el criterio: $500 > 125$ ms por factor 4
    \item Actualiza cada $2T_{coh}$ → pierde cambios del canal
    \item \textbf{NO cumple}
\end{itemize}

\textbf{Opción (c):} Pilotos cada 5 símbolos OFDM
\begin{itemize}
    \item Período: $T_c = 5 \times 250 = 1250$ ms
    \item Excede el criterio: $1250 > 125$ ms por factor 10
    \item Actualiza cada $5T_{coh}$ → inaceptable
    \item \textbf{NO cumple}
\end{itemize}

\textbf{Paso 5: Análisis del caso especial}

\textbf{Situación especial:} $T_{OFDM} = T_{coh} = 250$ ms

Cuando el símbolo OFDM dura exactamente el tiempo de coherencia:
\begin{itemize}
    \item El canal puede cambiar significativamente durante un símbolo
    \item El criterio estricto $T \leq T_{coh}/2$ requeriría símbolos más cortos
    \item Pero con las opciones disponibles, la (a) es la única aceptable
\end{itemize}

\textbf{Conclusión:}
\begin{itemize}
    \item Opción (a): Actualiza cada $T_{coh}$ - \textcolor{green}{cumple mínimamente}
    \item Opciones (b) y (c): Violan gravemente el criterio
\end{itemize}

\textbf{Respuesta:}

\[
\boxed{\text{Opción (a): Pilotos en todos los símbolos OFDM}}
\]

\textbf{Justificación detallada:}
\begin{itemize}
    \item \textbf{Cumple el criterio en caso límite:} Actualiza cada $T_{coh}$, que es aceptable cuando $T_{OFDM} = T_{coh}$
    \item \textbf{Es la única opción viable:} Las otras superan el criterio por factores de 4× y 10×
    \item \textbf{Minimiza el error de estimación:} Actualización continua del canal
    \item \textbf{Necesidad práctica:} En canales submarinos variables, perder una actualización compromete toda la comunicación
\end{itemize}

\textbf{¿Por qué no (b) o (c)?}
\begin{itemize}
    \item Esperarían 2 o 5 veces el tiempo de coherencia
    \item El canal cambiaría completamente entre actualizaciones
    \item La estimación sería obsoleta → alta tasa de error
    \item Pérdida total de sincronización con las variaciones del canal
\end{itemize}

\textbf{Compromiso necesario:}
\begin{itemize}
    \item Esta configuración usa pilotos en cada símbolo (overhead alto)
    \item Pero es \textbf{imprescindible} para seguir las variaciones del canal
    \item El débito se reduce, pero se gana \textbf{confiabilidad}
    \item Mejor 168 bits/s confiables que 256 bits/s con errores masivos
\end{itemize}

\textbf{Nota sobre el débito:} Esta configuración reduce el débito útil, pero es necesaria para la robustez del sistema en un canal que varía tan rápidamente.

\subsection*{Pregunta 4: Débito de información útil con pilotos}

\textbf{Enunciado:} Para la configuración de pilotos elegida, calcular el débito de información binaria útil del sistema.

\vspace{0.3cm}
\textbf{Solución paso a paso:}

\textbf{Paso 1: Configuración elegida}

De las preguntas anteriores:
\begin{itemize}
    \item Pilotos en frecuencia: 1 cada 3 subportadoras (11 pilotos)
    \item Pilotos en tiempo: en todos los símbolos OFDM
\end{itemize}

\textbf{Paso 2: Subportadoras de datos}

\[
N_{datos} = N - N_{pilotos} = 32 - 11 = 21 \text{ subportadoras}
\]

\textbf{Paso 3: Bits codificados por símbolo OFDM}

\[
\text{Bits codificados} = N_{datos} \times m = 21 \times 3 = 63 \text{ bits}
\]

\textbf{Paso 4: Bits de información (después del código corrector)}

Con rendimiento $r_c = \frac{2}{3}$:
\[
\text{Bits de información} = 63 \times \frac{2}{3} = 42 \text{ bits}
\]

\textbf{Paso 5: Débito binario útil}

Usamos el tiempo total de transmisión (incluyendo CP):
\[
T_{total} = 290 \text{ ms} = 0.29 \text{ s}
\]

\[
R_b = \frac{\text{Bits de información}}{T_{total}} = \frac{42}{0.29} \approx 144.83 \text{ bits/s}
\]

\textbf{Respuesta:}

\[
\boxed{R_b \approx 144.8 \text{ bits/s}}
\]

\textbf{Comparación:}
\begin{itemize}
    \item Sin pilotos: 220.7 bits/s
    \item Con pilotos: 144.8 bits/s
    \item Pérdida: $\frac{220.7-144.8}{220.7} \times 100\% \approx 34.4\%$
\end{itemize}

Este compromiso es necesario para poder estimar y ecualizar el canal correctamente.

\subsection*{Pregunta 5: Nueva configuración con $\Delta f = 8$ Hz}

\textbf{Enunciado:} Si ahora se elige el espaciamiento entre subportadoras $\Delta f = 8$ Hz, proponer un nuevo esquema de pilotos y calcular el nuevo débito útil.

\vspace{0.3cm}
\textbf{Solución paso a paso:}

\textbf{Paso 1: Calcular nueva $T_u$}

\[
T_u = \frac{1}{\Delta f} = \frac{1}{8} = 125 \text{ ms}
\]

\textbf{Paso 2: Calcular nuevo tiempo total}

Con el nuevo $\Delta f = 8$ Hz, el tiempo útil es:
\[
T_{OFDM,nuevo} = \frac{1}{\Delta f} = \frac{1}{8} = 125 \text{ ms}
\]

Manteniendo $T_{CP} = 40$ ms (el canal no ha cambiado):
\[
T_{total,nuevo} = T_{OFDM,nuevo} + T_{CP} = 125 + 40 = 165 \text{ ms}
\]

\textbf{Paso 3: Configuración de pilotos en frecuencia}

Criterio: $\Delta f_{piloto} \leq \frac{1}{2\tau_{max}} = 12.5$ Hz

Con $\Delta f = 8$ Hz:
\begin{itemize}
    \item 1 piloto cada 1 subportadora: $\Delta f_{p} = 8$ Hz $\leq 12.5$ Hz $\checkmark$
    \item Pero esto sería ineficiente (todas serían pilotos)
\end{itemize}

\textbf{Configuración óptima:} 1 piloto cada 1 subportadora (espaciamiento mínimo)

En realidad, podríamos usar:
\[
\text{Espaciamiento} = \lfloor 12.5 / 8 \rfloor = 1 \text{ subportadora}
\]

Propongo: \textbf{1 piloto cada 1 subportadora} o mejor dicho, debido a que 8 Hz ya cumple el criterio, podemos usar:

\textbf{Configuración propuesta:} 1 piloto cada 1 subportadora nos daría solo pilotos.

Mejor: Calculemos cuántas subportadoras entre pilotos:
\[
N_{spacing} = \lfloor \frac{12.5}{8} \rfloor = 1
\]

Esto indica que con $\Delta f = 8$ Hz, el criterio de Nyquist se satisface incluso con 1 piloto cada 2 subportadoras:
\[
\Delta f_{piloto} = 2 \times 8 = 16 \text{ Hz} > 12.5 \text{ Hz}
\]

Entonces, mejor configuración: \textbf{1 piloto cada 1 subportadora} (pero esto es extremo).

Propuesta práctica: \textbf{1 piloto cada 2 subportadoras}:
\begin{itemize}
    \item Número de pilotos: $32/2 = 16$ pilotos
    \item Subportadoras de datos: $32 - 16 = 16$
\end{itemize}

\textbf{Paso 4: Configuración de pilotos en tiempo}

Criterio: $T_{piloto} \leq \frac{T_{coh}}{2} = 125$ ms

Con $T_{OFDM} = 165$ ms $> 125$ ms, necesitamos pilotos en \textbf{todos los símbolos OFDM}.

\textbf{Paso 5: Calcular nuevo débito útil}

Con 16 subportadoras de datos y código $r_c = 2/3$:
\[
\text{Bits codificados/OFDM} = 16 \times 3 = 48 \text{ bits}
\]

\[
\text{Bits de información/OFDM} = 48 \times \frac{2}{3} = 32 \text{ bits}
\]

Usando el tiempo total:
\[
R_b = \frac{32}{0.165} \approx 193.94 \text{ bits/s}
\]

\textbf{Respuesta:}

\textbf{Esquema de pilotos propuesto:}
\[
\boxed{\begin{cases}
\text{Frecuencia: 1 piloto cada 2 subportadoras} \\
\text{Tiempo: pilotos en todos los símbolos OFDM}
\end{cases}}
\]

\textbf{Nuevo débito útil:}
\[
\boxed{R_b \approx 193.9 \text{ bits/s}}
\]

\textbf{Comparación con configuración anterior:}
\begin{itemize}
    \item Configuración 1 ($\Delta f = 4$ Hz, $T_{total}=290$ ms): 144.8 bits/s
    \item Configuración 2 ($\Delta f = 8$ Hz, $T_{total}=165$ ms): 193.9 bits/s
    \item Mejora: $\frac{193.9-144.8}{144.8} \times 100\% \approx 33.9\%$
\end{itemize}

\textbf{Ventajas de la nueva configuración:}
\begin{itemize}
    \item Mayor débito útil (+34\%)
    \item Menor latencia (165 ms vs 290 ms por símbolo)
    \item Mejor eficiencia del CP: $\frac{125}{165} = 75.8\%$ (vs 86.2\% anterior)
    \item Símbolos más cortos permiten mejor seguimiento del canal
\end{itemize}

\section*{Ejercicio 2 - Ecualización}

\textbf{Contexto del problema:}

Sistema de transmisión en banda base con:
\begin{itemize}
    \item Símbolos de información: $a_n \in \{-1, +1\}$
    \item Filtro de emisión: $h(t)$
    \item Canal de propagación: $c(t)$
    \item Ruido blanco aditivo gaussiano: $B(t)$
    \item Filtro de recepción: $g(t)$
    \item Muestras recibidas: $x_n = x(t_0 + nT)$
\end{itemize}

Modelo de la señal recibida:
\[
x(t_0 + nT) = A a_n + A a_{n-1} + b(t_0 + nT)
\]

donde $b(t) = (B * g)(t)$ y $A$ es un coeficiente positivo.

Regla de decisión:
\begin{itemize}
    \item Si $x(t_0 + nT) \geq 0 \Rightarrow \hat{a}_n = +1$
    \item Si $x(t_0 + nT) < 0 \Rightarrow \hat{a}_n = -1$
\end{itemize}

\subsection*{Pregunta 1: Probabilidad de error cuando $\sigma_b^2 \to 0$}

\textbf{Enunciado:} Sin hacer los cálculos pero justificando la respuesta, dar una estimación de la probabilidad de error cuando $\sigma_b^2$ (varianza del ruido) tiende a cero. \textcolor{red}{Atención: la probabilidad de error NO tiende a cero.}

\vspace{0.3cm}
\textbf{Solución:}

\textbf{Análisis del problema:}

\textit{Contexto:} Normalmente pensamos que sin ruido no habría errores. ¡Pero aquí hay una trampa! El problema es la \textbf{Interferencia Entre Símbolos (ISI)}.

\textit{¿Qué significa la ecuación?} Cuando recibimos la muestra en el instante $n$, no solo vemos el símbolo actual $a_n$, sino también una "huella" del símbolo anterior $a_{n-1}$:
\[
x_n = A a_n + A a_{n-1}
\]

\textit{Interpretación:} El canal tiene "memoria" - es como si tuviera eco. Lo que enviamos ahora se mezcla con lo que enviamos antes.

Vamos a analizar todos los casos posibles para ver cuándo esto causa problemas:

\textbf{Caso 1:} $a_n = +1$, $a_{n-1} = +1$
\[
x_n = A(+1) + A(+1) = +2A > 0 \Rightarrow \text{decisión correcta: } \hat{a}_n = +1 \quad \checkmark
\]

\textbf{Caso 2:} $a_n = +1$, $a_{n-1} = -1$
\[
x_n = A(+1) + A(-1) = 0 \Rightarrow \text{decisión: } \hat{a}_n = +1 \quad \checkmark
\]

\textbf{Caso 3:} $a_n = -1$, $a_{n-1} = +1$
\[
x_n = A(-1) + A(+1) = 0 \Rightarrow \text{decisión: } \hat{a}_n = +1 \quad \times
\]

\textcolor{red}{\textbf{¡Error!}} Debería decidir $\hat{a}_n = -1$ pero decide $+1$.

\textbf{Caso 4:} $a_n = -1$, $a_{n-1} = -1$
\[
x_n = A(-1) + A(-1) = -2A < 0 \Rightarrow \text{decisión correcta: } \hat{a}_n = -1 \quad \checkmark
\]

\textbf{Análisis de probabilidad:}

Suponiendo que los símbolos son equiprobables:
\begin{itemize}
    \item $P(a_n = +1) = P(a_n = -1) = \frac{1}{2}$
    \item $P(a_{n-1} = +1) = P(a_{n-1} = -1) = \frac{1}{2}$
\end{itemize}

La probabilidad de error ocurre solo en el Caso 3:
\[
P_e = P(a_n = -1 \text{ y } a_{n-1} = +1) = \frac{1}{2} \times \frac{1}{2} = \frac{1}{4}
\]

\textbf{Respuesta:}

\[
\boxed{P_e = \frac{1}{4} = 0.25 = 25\%}
\]

\textbf{Justificación:}

El error proviene de la \textbf{interferencia entre símbolos (ISI)}. Cuando el símbolo actual es $-1$ y el anterior es $+1$, ambas contribuciones se cancelan ($x_n = 0$), y por la regla de decisión ($x_n \geq 0 \Rightarrow \hat{a}_n = +1$), se decide incorrectamente.

Esta ISI no puede ser eliminada por reducir el ruido, de ahí que $P_e$ no tienda a cero incluso con SNR infinita.

\subsection*{Pregunta 2: Esquema de ecualización con realimentación de decisiones}

\textbf{Enunciado:} Para solucionar el problema, se desea realizar un ecualizador con realimentación de decisiones (DFE). Proponer el sinóptico correspondiente.

\vspace{0.3cm}
\textbf{Solución:}

\textbf{Principio del DFE:}

\textit{Idea central:} Si sabemos qué símbolos enviamos anteriormente, podemos \textbf{restar su contribución} de la señal recibida actual. Así eliminamos la interferencia del pasado.

\textit{¿Pero cómo sabemos los símbolos anteriores?} Usamos nuestras propias decisiones previas. Si decidimos que $\hat{a}_{n-1} = +1$, asumimos que eso es lo que se envió y restamos su efecto.

Un ecualizador DFE (Decision Feedback Equalizer) tiene dos componentes:

\begin{enumerate}
    \item \textbf{Filtro feedforward (FF):} Procesa la señal recibida. Compensa la ISI "precursora" (ecos que llegan antes del símbolo principal).
    
    \item \textbf{Filtro feedback (FB):} Usa las decisiones ya tomadas para cancelar la ISI "postcursora" (ecos que llegan después). \textit{Clave:} Este filtro resta la interferencia de símbolos pasados.
\end{enumerate}

\textit{Ventaja:} El filtro feedback no amplifica el ruido (solo usa decisiones binarias $\pm 1$).

\textbf{Esquema de bloques:}

\begin{center}
\begin{tikzpicture}[scale=0.9, transform shape]
    % Input
    \node [coordinate] (input) at (0,0) {};
    \node [above] at (input) {$x_n$};
    
    % Feedforward filter
    \node [draw, rectangle, minimum width=2.5cm, minimum height=1cm] (ff) at (2.5,0) {Filtro FF\\$\{c_k\}$};
    
    % Adder
    \node [draw, circle, minimum size=0.8cm] (adder) at (5.5,0) {$+$};
    
    % Decision device
    \node [draw, rectangle, minimum width=2cm, minimum height=1cm] (decision) at (8,0) {Decisión\\$\text{sign}(\cdot)$};
    
    % Output
    \node [coordinate] (output) at (10,0) {};
    \node [above] at (output) {$\hat{a}_n$};
    
    % Feedback filter
    \node [draw, rectangle, minimum width=2.5cm, minimum height=1cm] (fb) at (5.5,-2.5) {Filtro FB\\$\{b_k\}$};
    
    % Connections
    \draw[-latex] (input) -- (ff);
    \draw[-latex] (ff) -- node[above] {$y_n$} (adder);
    \draw[-latex] (adder) -- node[above] {$z_n$} (decision);
    \draw[-latex] (decision) -- (output);
    
    % Feedback connection
    \draw[-latex] (decision) |- (fb);
    \draw[-latex] (fb) -| node[pos=0.1, left] {$-$} (adder);
    
\end{tikzpicture}
\end{center}

\textbf{Funcionamiento detallado:}

\textbf{Filtro feedforward:} La señal recibida $x_n$ pasa por un filtro FIR con coeficientes $\{c_k\}$:
\[
y_n = \sum_{k=0}^{N_{FF}-1} c_k \cdot x_{n-k}
\]

\textbf{Filtro feedback:} Las decisiones previas $\hat{a}_{n-1}, \hat{a}_{n-2}, ...$ pasan por un filtro con coeficientes $\{b_k\}$:
\[
f_n = \sum_{k=1}^{N_{FB}} b_k \cdot \hat{a}_{n-k}
\]

\textbf{Combinación:} La señal de decisión es:
\[
z_n = y_n - f_n
\]

\textbf{Decisión:}
\[
\hat{a}_n = \text{sign}(z_n) = \begin{cases}
+1 & \text{si } z_n \geq 0 \\
-1 & \text{si } z_n < 0
\end{cases}
\]

\textbf{Para nuestro problema específico:}

Dado que tenemos ISI de un solo símbolo ($a_{n-1}$), necesitamos:
\begin{itemize}
    \item \textbf{Filtro FF:} Puede ser simple (incluso identidad $c_0 = 1$)
    \item \textbf{Filtro FB:} Un solo coeficiente $b_1$ para cancelar $A \cdot a_{n-1}$
\end{itemize}

Con $b_1 = A$, la señal de decisión sería:
\[
z_n = x_n - A \cdot \hat{a}_{n-1} = (A a_n + A a_{n-1}) - A \hat{a}_{n-1} = A a_n + A(a_{n-1} - \hat{a}_{n-1})
\]

Si $\hat{a}_{n-1} = a_{n-1}$ (decisión correcta anterior):
\[
z_n = A a_n
\]

¡La ISI ha sido cancelada!

\subsection*{Pregunta 3: Inconveniente del ecualizador DFE}

\textbf{Enunciado:} ¿Qué inconveniente presenta este tipo de ecualizador?

\vspace{0.3cm}
\textbf{Solución:}

\textbf{Inconveniente principal: Propagación de errores}

\textbf{Mecanismo:}

\begin{enumerate}
    \item Si se comete un error en la decisión de $\hat{a}_{n-1}$ (por ejemplo, debido al ruido)
    \item Este símbolo erróneo se usa en el filtro feedback para decidir $\hat{a}_n$
    \item La cancelación de ISI será incorrecta:
    \[
    z_n = A a_n + A(a_{n-1} - \hat{a}_{n-1})
    \]
    Si $\hat{a}_{n-1} \neq a_{n-1}$, tenemos un término de error $A(a_{n-1} - \hat{a}_{n-1})$ que puede ser hasta $\pm 2A$
    \item Este error aumenta la probabilidad de error en $\hat{a}_n$
    \item El error puede propagarse a símbolos futuros, creando \textbf{ráfagas de errores}
\end{enumerate}

\textbf{Ejemplo numérico:}

Supongamos $a_{n-1} = +1$ pero se decide incorrectamente $\hat{a}_{n-1} = -1$ (por ruido).

Para $a_n = +1$:
\[
x_n = A(+1) + A(+1) = 2A
\]
\[
z_n = 2A - A(-1) = 2A + A = 3A > 0 \quad \Rightarrow \quad \hat{a}_n = +1 \quad \checkmark
\]

Para $a_n = -1$:
\[
x_n = A(-1) + A(+1) = 0
\]
\[
z_n = 0 - A(-1) = A > 0 \quad \Rightarrow \quad \hat{a}_n = +1 \quad \times
\]

\textbf{Respuesta:}

\[
\boxed{\text{Propagación de errores (error propagation)}}
\]

\textbf{Otros inconvenientes secundarios:}
\begin{itemize}
    \item Degradación del rendimiento en canales con SNR baja
    \item Sensibilidad a errores de estimación del canal
    \item Dificultad para recuperarse de ráfagas de errores
\end{itemize}

\subsection*{Pregunta 4: Solución a la propagación de errores}

\textbf{Enunciado:} ¿Qué propone como medio para contrarrestar este inconveniente?

\vspace{0.3cm}
\textbf{Solución:}

\textbf{Varias soluciones posibles:}

\textbf{Solución 1: Secuencia de entrenamiento}

Transmitir periódicamente una secuencia conocida (training sequence) para:
\begin{itemize}
    \item Re-sincronizar el ecualizador
    \item Corregir errores acumulados
    \item Re-estimar los coeficientes del canal
\end{itemize}

\textbf{Implementación:}
\[
\text{Trama: } [\text{Training} | \text{Datos} | \text{Training} | \text{Datos} | ...]
\]

Durante la secuencia de entrenamiento, el feedback usa los símbolos conocidos en lugar de las decisiones.

\textbf{Solución 2: Código corrector de errores}

Usar un código corrector (por ejemplo, código convolucional o Reed-Solomon):
\begin{itemize}
    \item El código puede corregir errores aislados
    \item Reduce la probabilidad de que errores se propaguen
    \item Puede combinarse con entrelazado (interleaving) para dispersar errores en ráfaga
\end{itemize}

\textbf{Solución 3: Decisiones suaves con realimentación retardada}

En lugar de usar decisiones binarias inmediatamente:
\begin{itemize}
    \item Usar decisiones suaves (soft decisions) con información de fiabilidad
    \item Introducir un retardo para permitir decodificación con código corrector
    \item Realimentar solo decisiones con alta confianza
\end{itemize}

\textbf{Solución 4: Ecualización con algoritmo de Viterbi (MLSE)}

Cambiar a un ecualizador de máxima verosimilitud (ver pregunta 5):
\begin{itemize}
    \item Considera todas las secuencias posibles
    \item No sufre propagación de errores
    \item Más complejo computacionalmente
\end{itemize}

\textbf{Respuesta (solución práctica más común):}

\[
\boxed{\begin{array}{l}
\textbf{Uso de secuencias de entrenamiento periódicas} \\
\textbf{combinado con códigos correctores de errores}
\end{array}}
\]

\textbf{Justificación:}
\begin{itemize}
    \item Las secuencias de entrenamiento permiten "resetear" el ecualizador
    \item Los códigos correctores limitan la propagación de errores
    \item Compromiso razonable entre rendimiento y complejidad
    \item Solución estándar en sistemas prácticos (GSM, WiFi, etc.)
\end{itemize}

\subsection*{Pregunta 5: Ecualizador ML - Diagrama de Trellis}

\textbf{Enunciado:} En lugar de usar un ecualizador DFE, se quiere implementar un ecualizador ML (Maximum Likelihood). Dibujar el trellis correspondiente para bloques de tamaño 4. Se supone que el primer símbolo emitido es conocido y vale $-1$.

\vspace{0.3cm}
\textbf{Solución paso a paso:}

\textbf{Paso 1: Modelo del sistema como máquina de estados}

\textit{Nueva perspectiva:} En lugar de decidir símbolo por símbolo (como en DFE), vamos a pensar en \textbf{secuencias completas}. El ecualizador ML busca la secuencia más probable dado lo que recibimos.

\textit{¿Por qué estados?} Como el canal tiene memoria (cada símbolo depende del anterior), podemos modelarlo como una máquina de estados donde:
\[
x_n = A a_n + A a_{n-1} + b_n
\]

\textbf{Estado del sistema:} $S_n = a_{n-1}$ (el símbolo anterior)

\textit{¿Por qué es útil?} Conocer el estado nos dice qué esperar en la salida. Por ejemplo, si estamos en estado $S_n = +1$ y enviamos $a_n = +1$, sabemos que recibiremos aproximadamente $x_n \approx 2A$.

\textbf{Paso 2: Espacio de estados}

Estados posibles: $S \in \{-1, +1\}$ (2 estados)

\textbf{Paso 3: Transiciones}

Desde cada estado $S_n = a_{n-1}$, podemos emitir $a_n \in \{-1, +1\}$:

\textbf{Desde estado $S_n = -1$:}
\begin{itemize}
    \item Emitir $a_n = -1$: $x_n = A(-1) + A(-1) = -2A$, nuevo estado: $S_{n+1} = -1$
    \item Emitir $a_n = +1$: $x_n = A(+1) + A(-1) = 0$, nuevo estado: $S_{n+1} = +1$
\end{itemize}

\textbf{Desde estado $S_n = +1$:}
\begin{itemize}
    \item Emitir $a_n = -1$: $x_n = A(-1) + A(+1) = 0$, nuevo estado: $S_{n+1} = -1$
    \item Emitir $a_n = +1$: $x_n = A(+1) + A(+1) = +2A$, nuevo estado: $S_{n+1} = +1$
\end{itemize}

\textbf{Paso 4: Diagrama de Trellis para 4 símbolos}

Condición inicial: $a_0 = -1$ (conocido), por tanto $S_1 = -1$

Queremos decidir sobre: $a_1, a_2, a_3, a_4$ (4 símbolos)

\begin{center}
\begin{tikzpicture}[scale=0.85, transform shape]
    % Time labels
    \node at (-1, 1.5) {$t=0$};
    \node at (2, 1.5) {$t=1$};
    \node at (5, 1.5) {$t=2$};
    \node at (8, 1.5) {$t=3$};
    \node at (11, 1.5) {$t=4$};
    
    % States
    \node[draw, circle] (s0m1) at (0, 0) {$-1$};
    
    \node[draw, circle] (s1m1) at (3, -2) {$-1$};
    \node[draw, circle] (s1p1) at (3, 2) {$+1$};
    
    \node[draw, circle] (s2m1) at (6, -2) {$-1$};
    \node[draw, circle] (s2p1) at (6, 2) {$+1$};
    
    \node[draw, circle] (s3m1) at (9, -2) {$-1$};
    \node[draw, circle] (s3p1) at (9, 2) {$+1$};
    
    \node[draw, circle] (s4m1) at (12, -2) {$-1$};
    \node[draw, circle] (s4p1) at (12, 2) {$+1$};
    
    % Transitions t=0 to t=1
    \draw[-latex] (s0m1) -- (s1m1) node[midway, below, sloped] {\tiny $a_1=-1, x_1=-2A$};
    \draw[-latex] (s0m1) -- (s1p1) node[midway, above, sloped] {\tiny $a_1=+1, x_1=0$};
    
    % Transitions t=1 to t=2
    \draw[-latex] (s1m1) -- (s2m1) node[midway, below] {\tiny $a_2=-1, x_2=-2A$};
    \draw[-latex] (s1m1) -- (s2p1) node[midway, above, sloped] {\tiny $a_2=+1, x_2=0$};
    \draw[-latex] (s1p1) -- (s2m1) node[midway, below, sloped] {\tiny $a_2=-1, x_2=0$};
    \draw[-latex] (s1p1) -- (s2p1) node[midway, above] {\tiny $a_2=+1, x_2=+2A$};
    
    % Transitions t=2 to t=3
    \draw[-latex] (s2m1) -- (s3m1) node[midway, below] {\tiny $a_3=-1, x_3=-2A$};
    \draw[-latex] (s2m1) -- (s3p1) node[midway, above, sloped] {\tiny $a_3=+1, x_3=0$};
    \draw[-latex] (s2p1) -- (s3m1) node[midway, below, sloped] {\tiny $a_3=-1, x_3=0$};
    \draw[-latex] (s2p1) -- (s3p1) node[midway, above] {\tiny $a_3=+1, x_3=+2A$};
    
    % Transitions t=3 to t=4
    \draw[-latex] (s3m1) -- (s4m1) node[midway, below] {\tiny $a_4=-1, x_4=-2A$};
    \draw[-latex] (s3m1) -- (s4p1) node[midway, above, sloped] {\tiny $a_4=+1, x_4=0$};
    \draw[-latex] (s3p1) -- (s4m1) node[midway, below, sloped] {\tiny $a_4=-1, x_4=0$};
    \draw[-latex] (s3p1) -- (s4p1) node[midway, above] {\tiny $a_4=+1, x_4=+2A$};
    
\end{tikzpicture}
\end{center}

\textbf{Paso 5: Algoritmo de Viterbi - Búsqueda del mejor camino}

\textit{Idea genial:} Hay muchos caminos posibles a través del trellis ($2^4 = 16$ secuencias para 4 símbolos). ¿Cómo encontrar la mejor sin probar todas? \textbf{Viterbi lo hace eficientemente.}

\textbf{Métrica de rama:} Para cada transición (arco del trellis), medimos qué tan cerca está lo que recibimos de lo esperado:
\[
\mu_{rama} = |r_n - x_n|^2
\]

\textit{Interpretación:} Si recibimos $r_n$ y el camino predice $x_n$, el error cuadrático nos dice qué tan buena es esa transición.

\textbf{Métrica de camino acumulada:} Para cada estado, guardamos el mejor camino que llega hasta ahí:
\[
\mu_{camino}(S_n) = \min_{\text{todos los caminos}} \sum_{k=1}^{n} |r_k - x_k|^2
\]

\textit{Principio de optimalidad:} El mejor camino global debe pasar por los mejores caminos parciales. Por eso, en cada estado solo guardamos el camino "sobreviviente" con menor métrica.

\textbf{Procedimiento:}

\begin{enumerate}
    \item Inicializar: $\mu_0(-1) = 0$ (estado inicial conocido)
    \item Para cada instante $n = 1, 2, 3, 4$:
    \begin{itemize}
        \item Para cada estado $S_n$, calcular métrica de todas las ramas entrantes
        \item Seleccionar el camino con menor métrica (sobreviviente)
        \item Eliminar otros caminos (poda)
    \end{itemize}
    \item Al final, seleccionar el camino con menor métrica total
    \item Retroceder (traceback) para obtener la secuencia $\{a_1, a_2, a_3, a_4\}$
\end{enumerate}

\textbf{Ejemplo de cálculo:}

Supongamos recibimos: $r_1 = -1.8A$, $r_2 = 0.2A$, $r_3 = 1.9A$, $r_4 = 0.1A$

\textbf{En $t=1$:}

Camino 1: $a_1 = -1$, $x_1 = -2A$, métrica: $(-1.8A - (-2A))^2 = (0.2A)^2 = 0.04A^2$

Camino 2: $a_1 = +1$, $x_1 = 0$, métrica: $(-1.8A - 0)^2 = 3.24A^2$

$\Rightarrow$ Camino 1 es mejor (estado $-1$)

Y así sucesivamente para cada etapa...

\textbf{Ventajas del ecualizador ML:}

\begin{itemize}
    \item \textbf{Óptimo:} Minimiza la probabilidad de error de secuencia
    \item \textbf{No propaga errores:} Considera todos los caminos posibles
    \item \textbf{Mejor rendimiento:} Especialmente en canales severos
\end{itemize}

\textbf{Desventajas:}

\begin{itemize}
    \item \textbf{Complejidad:} Crece exponencialmente con la memoria del canal
    \item \textbf{Retardo:} Necesita procesar un bloque completo antes de decidir
    \item \textbf{Costo computacional:} Mayor que DFE
\end{itemize}

\textbf{Respuesta resumen:}

El trellis tiene 2 estados ($\pm 1$) y 4 etapas temporales. Desde cada estado parten 2 ramas (correspondientes a emitir $\pm 1$). El algoritmo de Viterbi selecciona el camino de menor métrica a través del trellis, proporcionando la secuencia más probable de símbolos transmitidos.

\vspace{1cm}

\section*{Resumen del examen}

Hemos resuelto completamente los 2 ejercicios del examen Final 2020:

\textbf{Ejercicio 1 - OFDM acústico submarino:}
\begin{itemize}
    \item Cálculo de parámetros del sistema: $\Delta f = 4$ Hz, $T_{CP} = 40$ ms, $T_{total} = 290$ ms
    \item Código corrector: $r_c = 2/3$
    \item Bits por símbolo OFDM: 64 bits (sin pilotos)
    \item Banda utilizada: 128 Hz alrededor de 12 kHz
    \item Débito útil sin pilotos: 220.7 bits/s
    \item Configuración óptima de pilotos: 1 cada 3 subportadoras (cumple Nyquist), en todos los símbolos
    \item Débito útil con pilotos: 144.8 bits/s
    \item Nueva configuración con $\Delta f = 8$ Hz: débito de 193.9 bits/s
\end{itemize}

\textbf{Ejercicio 2 - Ecualización:}
\begin{itemize}
    \item Probabilidad de error sin ruido (ISI): $P_e = 25\%$
    \item Diseño de ecualizador DFE con filtros feedforward y feedback
    \item Inconveniente principal: propagación de errores
    \item Soluciones: secuencias de entrenamiento y códigos correctores
    \item Ecualizador ML: diagrama de trellis con 2 estados y algoritmo de Viterbi
\end{itemize}

\end{document}
